% =============================================================
% results_template.tex
% Lead with the headline coefficient before walking through the
% table. Use \input{tabs/...} for every table.
% =============================================================

\section{Results}\label{sec:results}

\Cref{tab:main} presents our main results.

Column~(1) reports a baseline two-way fixed-effects specification with
no controls. The point estimate on [TREATMENT VARIABLE] is
\headlineEstimate{} (\headlineSE{}), statistically significant at the
[1/5] percent level.

Column~(2) adds [CONTROL SET]. The estimate moves to [b] ([se]) --
[stable / smaller / larger], suggesting that [INTERPRETATION].

Column~(3) adds [FIXED EFFECTS / ALTERNATIVE SPEC] and yields [b]
([se]). This is our preferred specification.

\input{tabs/table_main.tex}

To gauge economic significance, note that the headline effect of
\headlineEstimate{} corresponds to a [X]\% increase relative to the
sample mean of \headlineMean{} (\Cref{tab:summary}). Equivalently,
[ALTERNATIVE BENCHMARK].

\subsection{Heterogeneity}\label{sec:hetero}

\Cref{tab:hetero} explores heterogeneity by [DIMENSION]. The effect
is [LARGER / CONCENTRATED IN / DRIVEN BY] [SUBGROUP] (column~(2)),
with no detectable effect for [OTHER SUBGROUP] (column~(3)). This
pattern is consistent with [MECHANISM], which we examine more
directly in \Cref{sec:mechanisms}.

\input{tabs/table_hetero.tex}

\subsection{Dynamics}

\Cref{fig:event} reports the event-study coefficients from
\cref{eq:event_study}. Pre-period coefficients are economically and
statistically indistinguishable from zero, while post-period
coefficients trace out [SHAPE: a sharp jump / gradual ramp / inverted U].
The full dynamic ATT, aggregated across event time, is [b] ([se]).

\begin{figure}[htbp]\centering
  \includegraphics[width=0.85\textwidth]{figs/fig_event_study.pdf}
  \caption{Event-study estimates of the effect of [TREATMENT] on
  [OUTCOME]. Heterogeneity-robust Sun-Abraham (2021) estimator;
  95\% confidence intervals from cluster-robust standard errors
  clustered by [CLUSTER LEVEL].}
  \label{fig:event}
\end{figure}
